% Options for packages loaded elsewhere
\PassOptionsToPackage{unicode}{hyperref}
\PassOptionsToPackage{hyphens}{url}
\documentclass[
  11pt,
  a4paper,
]{article}
\usepackage{xcolor}
\usepackage[margin=18mm]{geometry}
\usepackage{amsmath,amssymb}
\setcounter{secnumdepth}{-\maxdimen} % remove section numbering
\usepackage{iftex}
\ifPDFTeX
  \usepackage[T1]{fontenc}
  \usepackage[utf8]{inputenc}
  \usepackage{textcomp} % provide euro and other symbols
\else % if luatex or xetex
  \usepackage{unicode-math} % this also loads fontspec
  \defaultfontfeatures{Scale=MatchLowercase}
  \defaultfontfeatures[\rmfamily]{Ligatures=TeX,Scale=1}
\fi
\ifPDFTeX\else
  % xetex/luatex font selection
    \setmainfont[]{Noto Serif}
    \setsansfont[]{Noto Sans}
    \setmonofont[]{Noto Sans Mono}
  \setmathfont[]{Noto Sans Math}
\fi
% Use upquote if available, for straight quotes in verbatim environments
\IfFileExists{upquote.sty}{\usepackage{upquote}}{}
\IfFileExists{microtype.sty}{% use microtype if available
  \usepackage[]{microtype}
  \UseMicrotypeSet[protrusion]{basicmath} % disable protrusion for tt fonts
}{}
\makeatletter
\@ifundefined{KOMAClassName}{% if non-KOMA class
  \IfFileExists{parskip.sty}{%
    \usepackage{parskip}
  }{% else
    \setlength{\parindent}{0pt}
    \setlength{\parskip}{6pt plus 2pt minus 1pt}}
}{% if KOMA class
  \KOMAoptions{parskip=half}}
\makeatother
\ifLuaTeX
\usepackage[bidi=basic]{babel}
\else
\usepackage[bidi=default]{babel}
\fi
\babelprovide[main,import]{russian}
\ifPDFTeX
\else
\babelfont{rm}[]{Noto Serif}
\fi
% get rid of language-specific shorthands (see #6817):
\let\LanguageShortHands\languageshorthands
\def\languageshorthands#1{}
\setlength{\emergencystretch}{3em} % prevent overfull lines
\providecommand{\tightlist}{%
  \setlength{\itemsep}{0pt}\setlength{\parskip}{0pt}}
\usepackage{bookmark}
\IfFileExists{xurl.sty}{\usepackage{xurl}}{} % add URL line breaks if available
\urlstyle{same}
\hypersetup{
  pdftitle={Билеты по дискретным случайным процессам},
  pdflang={ru-RU},
  hidelinks,
  pdfcreator={LaTeX via pandoc}}

\title{Билеты по дискретным случайным процессам}
\author{}
\date{}

\begin{document}
\maketitle

\section{Билет 05. Измеримые функции, отображения и случайные
элементы}\label{ux431ux438ux43bux435ux442-05.-ux438ux437ux43cux435ux440ux438ux43cux44bux435-ux444ux443ux43dux43aux446ux438ux438-ux43eux442ux43eux431ux440ux430ux436ux435ux43dux438ux44f-ux438-ux441ux43bux443ux447ux430ux439ux43dux44bux435-ux44dux43bux435ux43cux435ux43dux442ux44b}

Источники: Ширяев 1, гл. II; Сердобольская, лекция 1; лекции ДСП.

\subsection{1. Короткий
ответ}\label{ux43aux43eux440ux43eux442ux43aux438ux439-ux43eux442ux432ux435ux442}

Измеримое отображение - это отображение между измеримыми пространствами,
у которого прообраз любого измеримого множества снова измерим. Если

\[
X:(\Omega,\mathcal F)\to(E,\mathcal E),
\]

то измеримость означает

\[
X^{-1}(B)=\{\omega\in\Omega:X(\omega)\in B\}\in\mathcal F
\qquad \forall B\in\mathcal E.
\]

Случайная величина - это измеримое отображение из вероятностного
пространства в \(\mathbb R\) с борелевской \(\sigma\)-алгеброй:

\[
X:(\Omega,\mathcal F,P)\to(\mathbb R,\mathcal B(\mathbb R)).
\]

Случайный вектор - измеримое отображение в \(\mathbb R^d\), случайный
элемент - измеримое отображение в произвольное измеримое, обычно
метрическое, пространство. Измеримость нужна для того, чтобы события
вида \(\{X\in B\}\) имели вероятность.

\subsection{2. Полный
ответ}\label{ux43fux43eux43bux43dux44bux439-ux43eux442ux432ux435ux442}

\subsubsection{2.1. Мотивация и связь с наблюдаемыми
величинами}\label{ux43cux43eux442ux438ux432ux430ux446ux438ux44f-ux438-ux441ux432ux44fux437ux44c-ux441-ux43dux430ux431ux43bux44eux434ux430ux435ux43cux44bux43cux438-ux432ux435ux43bux438ux447ux438ux43dux430ux43cux438}

В предыдущих билетах вводились \(\sigma\)-алгебры, меры, вероятностные
пространства и борелевские множества. Теперь нужно связать пространство
элементарных исходов с наблюдаемыми величинами. В вероятности мы редко
работаем напрямую с исходом \(\omega\in\Omega\). Обычно нас интересует
численная или более сложная характеристика исхода: результат измерения,
вектор наблюдений, траектория процесса, последовательность значений.

\subsubsection{2.2. Определение измеримого
отображения}\label{ux43eux43fux440ux435ux434ux435ux43bux435ux43dux438ux435-ux438ux437ux43cux435ux440ux438ux43cux43eux433ux43e-ux43eux442ux43eux431ux440ux430ux436ux435ux43dux438ux44f}

Пусть есть измеримое пространство \((\Omega,\mathcal F)\) и другое
измеримое пространство \((E,\mathcal E)\). Отображение

\[
X:\Omega\to E
\]

называется измеримым, если для любого измеримого множества
\(B\in\mathcal E\) его прообраз измерим:

\[
X^{-1}(B)\in\mathcal F.
\]

Интуитивно это означает: если в пространстве значений \(E\) можно задать
наблюдаемое событие \(B\), то событие ``значение \(X\) попало в \(B\)''
должно быть событием в исходном вероятностном пространстве. Иначе
выражение \(P\{X\in B\}\) просто не будет определено.

\subsubsection{2.3. Вещественная случайная
величина}\label{ux432ux435ux449ux435ux441ux442ux432ux435ux43dux43dux430ux44f-ux441ux43bux443ux447ux430ux439ux43dux430ux44f-ux432ux435ux43bux438ux447ux438ux43dux430}

Главный частный случай - вещественная случайная величина. Это измеримое
отображение

\[
X:(\Omega,\mathcal F,P)\to(\mathbb R,\mathcal B(\mathbb R)).
\]

По определению для любого борелевского множества \(B\subset\mathbb R\)
должно быть

\[
\{X\in B\}\in\mathcal F.
\]

Проверять все борелевские множества напрямую неудобно. Поэтому
используют порождающие классы. Так как

\[
\mathcal B(\mathbb R)=\sigma\{(-\infty,x]:x\in\mathbb R\},
\]

то для вещественной функции достаточно проверить, что

\[
\{X\le x\}\in\mathcal F\qquad \forall x\in\mathbb R.
\]

Эквивалентно можно проверять множества \(\{X<x\}\), \(\{X>x\}\) или
\(\{X\ge x\}\) для всех \(x\).

\subsubsection{2.4. Случайный вектор и критерий
измеримости}\label{ux441ux43bux443ux447ux430ux439ux43dux44bux439-ux432ux435ux43aux442ux43eux440-ux438-ux43aux440ux438ux442ux435ux440ux438ux439-ux438ux437ux43cux435ux440ux438ux43cux43eux441ux442ux438}

Случайный вектор

\[
X=(X_1,\ldots,X_d)
\]

со значениями в \(\mathbb R^d\) - это измеримое отображение

\[
X:(\Omega,\mathcal F,P)\to(\mathbb R^d,\mathcal B(\mathbb R^d)).
\]

Важный критерий: \(X\) измерим тогда и только тогда, когда все его
координаты \(X_1,\ldots,X_d\) измеримы как вещественные случайные
величины. Это удобно, потому что на практике вектор обычно задается
координатами.

\subsubsection{2.5. Случайный элемент в метрическом
пространстве}\label{ux441ux43bux443ux447ux430ux439ux43dux44bux439-ux44dux43bux435ux43cux435ux43dux442-ux432-ux43cux435ux442ux440ux438ux447ux435ux441ux43aux43eux43c-ux43fux440ux43eux441ux442ux440ux430ux43dux441ux442ux432ux435}

Еще более общий объект - случайный элемент. Если \((E,\rho)\) -
метрическое пространство, то случайным элементом со значениями в \(E\)
называют измеримое отображение

\[
X:(\Omega,\mathcal F,P)\to(E,\mathcal B(E)),
\]

где \(\mathcal B(E)\) - борелевская \(\sigma\)-алгебра метрического
пространства. Примеры: случайная точка в \(\mathbb R^d\), случайная
последовательность в \(\mathbb R^{\mathbb N}\), случайная функция в
пространстве траекторий.

\subsubsection{2.6. Измеримость случайных
последовательностей}\label{ux438ux437ux43cux435ux440ux438ux43cux43eux441ux442ux44c-ux441ux43bux443ux447ux430ux439ux43dux44bux445-ux43fux43eux441ux43bux435ux434ux43eux432ux430ux442ux435ux43bux44cux43dux43eux441ux442ux435ux439}

Для дискретных случайных процессов это особенно важно.
Последовательность случайных величин \((X_1,X_2,\ldots)\) можно
рассматривать двояко:

\begin{enumerate}
\def\labelenumi{\arabic{enumi}.}
\tightlist
\item
  как семейство измеримых отображений \(X_n:\Omega\to\mathbb R\);
\item
  как одно отображение
\end{enumerate}

\[
X:\Omega\to\mathbb R^{\mathbb N},\qquad
X(\omega)=(X_1(\omega),X_2(\omega),\ldots).
\]

Если на \(\mathbb R^{\mathbb N}\) взять цилиндрическую, то есть
произведенную, \(\sigma\)-алгебру, то такое отображение измеримо тогда и
только тогда, когда измеримы все координаты \(X_n\). Это будет
использоваться в билетах про конечномерные распределения, теорему
Колмогорова и дискретные случайные процессы.

\subsubsection{2.7. Роль условия измеримости в теории
вероятностей}\label{ux440ux43eux43bux44c-ux443ux441ux43bux43eux432ux438ux44f-ux438ux437ux43cux435ux440ux438ux43cux43eux441ux442ux438-ux432-ux442ux435ux43eux440ux438ux438-ux432ux435ux440ux43eux44fux442ux43dux43eux441ux442ux435ux439}

Итак, роль билета такая: измеримость - минимальное техническое условие,
которое превращает функцию от исхода в корректный вероятностный объект.
После этого можно задавать распределение \(P_X(B)=P\{X\in B\}\),
интегрировать функции от \(X\), говорить о моментах, независимости,
условном ожидании и процессах.

\subsection{3. Основные
определения}\label{ux43eux441ux43dux43eux432ux43dux44bux435-ux43eux43fux440ux435ux434ux435ux43bux435ux43dux438ux44f}

\subsubsection{Измеримое
пространство}\label{ux438ux437ux43cux435ux440ux438ux43cux43eux435-ux43fux440ux43eux441ux442ux440ux430ux43dux441ux442ux432ux43e}

Пара \((\Omega,\mathcal F)\), где \(\Omega\) - множество, а
\(\mathcal F\) - \(\sigma\)-алгебра его подмножеств.

\subsubsection{Прообраз
множества}\label{ux43fux440ux43eux43eux431ux440ux430ux437-ux43cux43dux43eux436ux435ux441ux442ux432ux430}

Если \(X:\Omega\to E\) и \(B\subset E\), то

\[
X^{-1}(B)=\{\omega\in\Omega:X(\omega)\in B\}.
\]

Важно: прообраз всегда является подмножеством \(\Omega\), а не
подмножеством \(E\).

\subsubsection{Измеримое
отображение}\label{ux438ux437ux43cux435ux440ux438ux43cux43eux435-ux43eux442ux43eux431ux440ux430ux436ux435ux43dux438ux435}

Пусть \((\Omega,\mathcal F)\) и \((E,\mathcal E)\) - измеримые
пространства. Отображение

\[
X:\Omega\to E
\]

называется \(\mathcal F/\mathcal E\)-измеримым, если

\[
X^{-1}(B)\in\mathcal F
\qquad \forall B\in\mathcal E.
\]

\subsubsection{Борелевски измеримая
функция}\label{ux431ux43eux440ux435ux43bux435ux432ux441ux43aux438-ux438ux437ux43cux435ux440ux438ux43cux430ux44f-ux444ux443ux43dux43aux446ux438ux44f}

Если \(E\) - топологическое или метрическое пространство и
\(\mathcal E=\mathcal B(E)\), то измеримое отображение в
\((E,\mathcal B(E))\) называют борелевски измеримым.

\subsubsection{Вещественная случайная
величина}\label{ux432ux435ux449ux435ux441ux442ux432ux435ux43dux43dux430ux44f-ux441ux43bux443ux447ux430ux439ux43dux430ux44f-ux432ux435ux43bux438ux447ux438ux43dux430-1}

Измеримое отображение

\[
X:(\Omega,\mathcal F,P)\to(\mathbb R,\mathcal B(\mathbb R)).
\]

То есть для любого \(B\in\mathcal B(\mathbb R)\) событие \(\{X\in B\}\)
принадлежит \(\mathcal F\).

\subsubsection{Случайный
вектор}\label{ux441ux43bux443ux447ux430ux439ux43dux44bux439-ux432ux435ux43aux442ux43eux440}

Измеримое отображение

\[
X=(X_1,\ldots,X_d):(\Omega,\mathcal F,P)\to(\mathbb R^d,\mathcal B(\mathbb R^d)).
\]

\subsubsection{Случайный элемент метрического
пространства}\label{ux441ux43bux443ux447ux430ux439ux43dux44bux439-ux44dux43bux435ux43cux435ux43dux442-ux43cux435ux442ux440ux438ux447ux435ux441ux43aux43eux433ux43e-ux43fux440ux43eux441ux442ux440ux430ux43dux441ux442ux432ux430}

Если \((E,\rho)\) - метрическое пространство, то случайный элемент со
значениями в \(E\) - это измеримое отображение

\[
X:(\Omega,\mathcal F,P)\to(E,\mathcal B(E)).
\]

\subsubsection{Случайная
последовательность}\label{ux441ux43bux443ux447ux430ux439ux43dux430ux44f-ux43fux43eux441ux43bux435ux434ux43eux432ux430ux442ux435ux43bux44cux43dux43eux441ux442ux44c}

Случайный элемент пространства последовательностей:

\[
X:\Omega\to\mathbb R^{\mathbb N},\qquad
X(\omega)=(X_1(\omega),X_2(\omega),\ldots),
\]

обычно с произведенной \(\sigma\)-алгеброй

\[
\mathcal B(\mathbb R)^{\otimes\mathbb N}
=
\sigma\{\pi_n^{-1}(B):n\in\mathbb N,\ B\in\mathcal B(\mathbb R)\},
\]

где \(\pi_n(x_1,x_2,\ldots)=x_n\) - координатная проекция.

\subsubsection{Сечение
процесса}\label{ux441ux435ux447ux435ux43dux438ux435-ux43fux440ux43eux446ux435ux441ux441ux430}

Для случайного процесса \(\{X_t:t\in T\}\) фиксированное значение
\(X_t(\omega)\) как функция \(\omega\) называется сечением процесса в
момент \(t\).

\subsubsection{Траектория
процесса}\label{ux442ux440ux430ux435ux43aux442ux43eux440ux438ux44f-ux43fux440ux43eux446ux435ux441ux441ux430}

Для фиксированного \(\omega\) функция

\[
t\mapsto X_t(\omega)
\]

называется траекторией процесса.

\subsection{4. Основные теоремы и
свойства}\label{ux43eux441ux43dux43eux432ux43dux44bux435-ux442ux435ux43eux440ux435ux43cux44b-ux438-ux441ux432ux43eux439ux441ux442ux432ux430}

\subsubsection{Свойство 1. Достаточно проверять порождающий
класс}\label{ux441ux432ux43eux439ux441ux442ux432ux43e-1.-ux434ux43eux441ux442ux430ux442ux43eux447ux43dux43e-ux43fux440ux43eux432ux435ux440ux44fux442ux44c-ux43fux43eux440ux43eux436ux434ux430ux44eux449ux438ux439-ux43aux43bux430ux441ux441}

Пусть \((\Omega,\mathcal F)\) и \((E,\mathcal E)\) - измеримые
пространства, а \(\mathcal E=\sigma(\mathcal C)\). Тогда отображение
\(X:\Omega\to E\) измеримо тогда и только тогда, когда

\[
X^{-1}(C)\in\mathcal F
\qquad \forall C\in\mathcal C.
\]

\subsubsection{Доказательство}\label{ux434ux43eux43aux430ux437ux430ux442ux435ux43bux44cux441ux442ux432ux43e}

Рассмотрим класс

\[
\mathcal D=\{B\subset E:X^{-1}(B)\in\mathcal F\}.
\]

Этот класс является \(\sigma\)-алгеброй. Действительно,

\[
X^{-1}(E)=\Omega\in\mathcal F,
\]

\[
X^{-1}(E\setminus B)=\Omega\setminus X^{-1}(B),
\]

и

\[
X^{-1}\left(\bigcup_{n=1}^{\infty}B_n\right)
=
\bigcup_{n=1}^{\infty}X^{-1}(B_n).
\]

Если \(X^{-1}(C)\in\mathcal F\) для всех \(C\in\mathcal C\), то
\(\mathcal C\subset\mathcal D\). Так как \(\mathcal D\) -
\(\sigma\)-алгебра, она содержит \(\sigma(\mathcal C)=\mathcal E\).
Значит \(X\) измеримо. Обратное направление следует прямо из
определения. \(\square\)

\subsubsection{Свойство 2. Критерии измеримости вещественной
функции}\label{ux441ux432ux43eux439ux441ux442ux432ux43e-2.-ux43aux440ux438ux442ux435ux440ux438ux438-ux438ux437ux43cux435ux440ux438ux43cux43eux441ux442ux438-ux432ux435ux449ux435ux441ux442ux432ux435ux43dux43dux43eux439-ux444ux443ux43dux43aux446ux438ux438}

Для функции \(X:\Omega\to\mathbb R\) следующие условия эквивалентны:

\begin{enumerate}
\def\labelenumi{\arabic{enumi}.}
\tightlist
\item
  \(X\) измерима как отображение в
  \((\mathbb R,\mathcal B(\mathbb R))\);
\item
  \(\{X\le x\}\in\mathcal F\) для всех \(x\in\mathbb R\);
\item
  \(\{X<x\}\in\mathcal F\) для всех \(x\in\mathbb R\);
\item
  \(\{X>x\}\in\mathcal F\) для всех \(x\in\mathbb R\);
\item
  \(\{X\ge x\}\in\mathcal F\) для всех \(x\in\mathbb R\).
\end{enumerate}

\subsubsection{Доказательство-идея}\label{ux434ux43eux43aux430ux437ux430ux442ux435ux43bux44cux441ux442ux432ux43e-ux438ux434ux435ux44f}

Борелевская \(\sigma\)-алгебра на прямой порождается лучами:

\[
\mathcal B(\mathbb R)=\sigma\{(-\infty,x]:x\in\mathbb R\}.
\]

По свойству 1 достаточно проверять прообразы таких лучей:

\[
X^{-1}((-\infty,x])=\{X\le x\}.
\]

Остальные варианты получаются через дополнения и счетные пересечения:

\[
\{X<x\}=\bigcup_{n=1}^{\infty}\{X\le x-\tfrac1n\},
\]

\[
\{X\ge x\}=\Omega\setminus\{X<x\}.
\]

\subsubsection{Свойство 3. Для борелевской измеримости достаточно
открытых
множеств}\label{ux441ux432ux43eux439ux441ux442ux432ux43e-3.-ux434ux43bux44f-ux431ux43eux440ux435ux43bux435ux432ux441ux43aux43eux439-ux438ux437ux43cux435ux440ux438ux43cux43eux441ux442ux438-ux434ux43eux441ux442ux430ux442ux43eux447ux43dux43e-ux43eux442ux43aux440ux44bux442ux44bux445-ux43cux43dux43eux436ux435ux441ux442ux432}

Пусть \(E\) - топологическое пространство. Отображение

\[
X:(\Omega,\mathcal F)\to(E,\mathcal B(E))
\]

измеримо тогда и только тогда, когда

\[
X^{-1}(G)\in\mathcal F
\]

для любого открытого множества \(G\subset E\).

\subsubsection{Доказательство}\label{ux434ux43eux43aux430ux437ux430ux442ux435ux43bux44cux441ux442ux432ux43e-1}

По определению

\[
\mathcal B(E)=\sigma\{G\subset E:G\text{ открыто}\}.
\]

Остается применить свойство 1.

\subsubsection{Свойство 4. Непрерывные отображения борелевски
измеримы}\label{ux441ux432ux43eux439ux441ux442ux432ux43e-4.-ux43dux435ux43fux440ux435ux440ux44bux432ux43dux44bux435-ux43eux442ux43eux431ux440ux430ux436ux435ux43dux438ux44f-ux431ux43eux440ux435ux43bux435ux432ux441ux43aux438-ux438ux437ux43cux435ux440ux438ux43cux44b}

Пусть \(E\) и \(S\) - топологические пространства, а

\[
f:E\to S
\]

непрерывно. Тогда \(f\) измеримо как отображение

\[
(E,\mathcal B(E))\to(S,\mathcal B(S)).
\]

\subsubsection{Доказательство}\label{ux434ux43eux43aux430ux437ux430ux442ux435ux43bux44cux441ux442ux432ux43e-2}

Для любого открытого \(G\subset S\) прообраз \(f^{-1}(G)\) открыт в
\(E\), значит принадлежит \(\mathcal B(E)\). По свойству 3 отображение
\(f\) борелевски измеримо. \(\square\)

\subsubsection{Свойство 5. Композиция измеримых отображений
измерима}\label{ux441ux432ux43eux439ux441ux442ux432ux43e-5.-ux43aux43eux43cux43fux43eux437ux438ux446ux438ux44f-ux438ux437ux43cux435ux440ux438ux43cux44bux445-ux43eux442ux43eux431ux440ux430ux436ux435ux43dux438ux439-ux438ux437ux43cux435ux440ux438ux43cux430}

Если

\[
X:(\Omega,\mathcal F)\to(E,\mathcal E)
\]

измеримо и

\[
f:(E,\mathcal E)\to(S,\mathcal S)
\]

измеримо, то композиция

\[
f\circ X:(\Omega,\mathcal F)\to(S,\mathcal S)
\]

измерима.

\subsubsection{Доказательство}\label{ux434ux43eux43aux430ux437ux430ux442ux435ux43bux44cux441ux442ux432ux43e-3}

Для любого \(C\in\mathcal S\)

\[
(f\circ X)^{-1}(C)=X^{-1}(f^{-1}(C)).
\]

Так как \(f\) измеримо, \(f^{-1}(C)\in\mathcal E\). Так как \(X\)
измеримо, \(X^{-1}(f^{-1}(C))\in\mathcal F\). \(\square\)

\subsubsection{Свойство 6. Критерий измеримости случайного
вектора}\label{ux441ux432ux43eux439ux441ux442ux432ux43e-6.-ux43aux440ux438ux442ux435ux440ux438ux439-ux438ux437ux43cux435ux440ux438ux43cux43eux441ux442ux438-ux441ux43bux443ux447ux430ux439ux43dux43eux433ux43e-ux432ux435ux43aux442ux43eux440ux430}

Отображение

\[
X=(X_1,\ldots,X_d):\Omega\to\mathbb R^d
\]

измеримо относительно \(\mathcal B(\mathbb R^d)\) тогда и только тогда,
когда все координаты

\[
X_1,\ldots,X_d
\]

измеримы как вещественные функции.

\subsubsection{Доказательство}\label{ux434ux43eux43aux430ux437ux430ux442ux435ux43bux44cux441ux442ux432ux43e-4}

Если \(X\) измеримо, то координата

\[
X_k=\pi_k\circ X
\]

измерима, потому что проекция \(\pi_k:\mathbb R^d\to\mathbb R\)
непрерывна, а композиция измеримых отображений измерима.

Обратно, пусть все \(X_k\) измеримы. Достаточно проверить прообразы
прямоугольников

\[
(-\infty,x_1]\times\cdots\times(-\infty,x_d],
\]

потому что такие множества порождают \(\mathcal B(\mathbb R^d)\). Но

\[
X^{-1}\left((-\infty,x_1]\times\cdots\times(-\infty,x_d]\right)
=
\bigcap_{k=1}^d\{X_k\le x_k\}.
\]

Каждое множество справа лежит в \(\mathcal F\), а \(\mathcal F\)
замкнута относительно конечных пересечений. Значит \(X\) измерим.
\(\square\)

\subsubsection{Свойство 7. Борелевская функция от случайного вектора
дает случайную
величину}\label{ux441ux432ux43eux439ux441ux442ux432ux43e-7.-ux431ux43eux440ux435ux43bux435ux432ux441ux43aux430ux44f-ux444ux443ux43dux43aux446ux438ux44f-ux43eux442-ux441ux43bux443ux447ux430ux439ux43dux43eux433ux43e-ux432ux435ux43aux442ux43eux440ux430-ux434ux430ux435ux442-ux441ux43bux443ux447ux430ux439ux43dux443ux44e-ux432ux435ux43bux438ux447ux438ux43dux443}

Если \(X:\Omega\to\mathbb R^d\) - случайный вектор, а
\(g:\mathbb R^d\to\mathbb R^m\) - борелевски измеримое отображение, то

\[
g(X)
\]

является случайным вектором в \(\mathbb R^m\).

В частности, если \(X,Y\) - вещественные случайные величины, то
измеримы:

\[
X+Y,\qquad XY,\qquad \max(X,Y),\qquad \min(X,Y),
\]

если \(Y\ne0\) всюду, то также \(X/Y\). Если \(Y=0\) возможно, нужно
определить значение на множестве \(\{Y=0\}\), например положить там
\(0\), и рассматривать такую измеримую модификацию.

\subsubsection{Обоснование}\label{ux43eux431ux43eux441ux43dux43eux432ux430ux43dux438ux435}

Вектор \((X,Y)\) измерим по свойству 6, а операции сложения, умножения,
максимума и минимума являются непрерывными функциями
\(\mathbb R^2\to\mathbb R\).

\subsubsection{Свойство 8. Предел измеримых вещественных функций
измерим}\label{ux441ux432ux43eux439ux441ux442ux432ux43e-8.-ux43fux440ux435ux434ux435ux43b-ux438ux437ux43cux435ux440ux438ux43cux44bux445-ux432ux435ux449ux435ux441ux442ux432ux435ux43dux43dux44bux445-ux444ux443ux43dux43aux446ux438ux439-ux438ux437ux43cux435ux440ux438ux43c}

Если \(X_n:\Omega\to\mathbb R\) измеримы, то функции

\[
\sup_n X_n,\qquad \inf_n X_n,\qquad \limsup_{n\to\infty}X_n,\qquad \liminf_{n\to\infty}X_n
\]

измеримы как расширенные вещественные функции. Если
\(X_n(\omega)\to X(\omega)\) для всех \(\omega\), то \(X\) измерима.

\subsubsection{Доказательство-идея}\label{ux434ux43eux43aux430ux437ux430ux442ux435ux43bux44cux441ux442ux432ux43e-ux438ux434ux435ux44f-1}

Например,

\[
\left\{\sup_n X_n>a\right\}
=
\bigcup_{n=1}^{\infty}\{X_n>a\}\in\mathcal F.
\]

Значит \(\sup_nX_n\) измерима. Инфимум выражается через супремум:

\[
\inf_nX_n=-\sup_n(-X_n).
\]

Далее

\[
\limsup_{n\to\infty}X_n=\inf_{m\ge1}\sup_{n\ge m}X_n,
\]

\[
\liminf_{n\to\infty}X_n=\sup_{m\ge1}\inf_{n\ge m}X_n.
\]

Если обычный предел существует, то

\[
X=\limsup_{n\to\infty}X_n=\liminf_{n\to\infty}X_n.
\]

\subsubsection{Свойство 9. Измеримость случайной
последовательности}\label{ux441ux432ux43eux439ux441ux442ux432ux43e-9.-ux438ux437ux43cux435ux440ux438ux43cux43eux441ux442ux44c-ux441ux43bux443ux447ux430ux439ux43dux43eux439-ux43fux43eux441ux43bux435ux434ux43eux432ux430ux442ux435ux43bux44cux43dux43eux441ux442ux438}

Пусть

\[
X(\omega)=(X_1(\omega),X_2(\omega),\ldots)
\]

и на \(\mathbb R^{\mathbb N}\) взята произведенная \(\sigma\)-алгебра
\(\mathcal B(\mathbb R)^{\otimes\mathbb N}\). Тогда \(X\) измеримо тогда
и только тогда, когда каждая координата \(X_n\) измерима.

\subsubsection{Доказательство-идея}\label{ux434ux43eux43aux430ux437ux430ux442ux435ux43bux44cux441ux442ux432ux43e-ux438ux434ux435ux44f-2}

Произведенная \(\sigma\)-алгебра порождается цилиндрами

\[
\pi_{i_1}^{-1}(B_1)\cap\cdots\cap\pi_{i_k}^{-1}(B_k),
\]

где \(B_j\in\mathcal B(\mathbb R)\). Прообраз такого цилиндра равен

\[
\{X_{i_1}\in B_1\}\cap\cdots\cap\{X_{i_k}\in B_k\},
\]

что измеримо при измеримости всех координат.

\subsection{5. Примеры}\label{ux43fux440ux438ux43cux435ux440ux44b}

\subsubsection{Пример 1. Индикатор
события}\label{ux43fux440ux438ux43cux435ux440-1.-ux438ux43dux434ux438ux43aux430ux442ux43eux440-ux441ux43eux431ux44bux442ux438ux44f}

Пусть \(A\in\mathcal F\). Тогда

\[
\mathbf 1_A(\omega)=
\begin{cases}
1, & \omega\in A,\\
0, & \omega\notin A
\end{cases}
\]

является случайной величиной. Действительно, для любого \(x\) множество
\(\{\mathbf 1_A\le x\}\) равно одному из множеств

\[
\varnothing,\qquad \Omega\setminus A,\qquad \Omega,
\]

и все они принадлежат \(\mathcal F\).

\subsubsection{Пример 2. Простая случайная
величина}\label{ux43fux440ux438ux43cux435ux440-2.-ux43fux440ux43eux441ux442ux430ux44f-ux441ux43bux443ux447ux430ux439ux43dux430ux44f-ux432ux435ux43bux438ux447ux438ux43dux430}

Если \(A_1,\ldots,A_n\in\mathcal F\) и \(a_1,\ldots,a_n\in\mathbb R\),
то

\[
X(\omega)=\sum_{k=1}^n a_k\mathbf 1_{A_k}(\omega)
\]

измерима. Она принимает конечное число значений. Чтобы увидеть
измеримость строго, можно разбить \(\Omega\) на конечное число атомов
вида

\[
\bigcap_{k=1}^n C_k,\qquad C_k\in\{A_k,\Omega\setminus A_k\}.
\]

На каждом таком атоме \(X\) постоянна, поэтому \(\{X\in B\}\) является
конечным объединением измеримых атомов. Такие функции будут основой
конструкции интеграла Лебега в Билете 07.

\subsubsection{Пример 3. Случайный
вектор}\label{ux43fux440ux438ux43cux435ux440-3.-ux441ux43bux443ux447ux430ux439ux43dux44bux439-ux432ux435ux43aux442ux43eux440}

Если \(X\) и \(Y\) - случайные величины, то

\[
Z=(X,Y):\Omega\to\mathbb R^2
\]

является случайным вектором. Например, событие попадания в прямоугольник
имеет вид

\[
\{Z\in(-\infty,x]\times(-\infty,y]\}
=
\{X\le x\}\cap\{Y\le y\}.
\]

Это событие измеримо.

\subsubsection{Пример 4. Случайная
последовательность}\label{ux43fux440ux438ux43cux435ux440-4.-ux441ux43bux443ux447ux430ux439ux43dux430ux44f-ux43fux43eux441ux43bux435ux434ux43eux432ux430ux442ux435ux43bux44cux43dux43eux441ux442ux44c}

Пусть \(X_1,X_2,\ldots\) - случайные величины. Тогда

\[
X(\omega)=(X_1(\omega),X_2(\omega),\ldots)
\]

является случайным элементом пространства \(\mathbb R^{\mathbb N}\) с
произведенной \(\sigma\)-алгеброй. События, зависящие от конечного числа
координат, например

\[
\{X_1\le a_1,\ X_2\le a_2,\ X_5\in B\},
\]

измеримы. Именно такие события задают конечномерные распределения.

\subsubsection{Пример 5. Непрерывная функция от случайной
величины}\label{ux43fux440ux438ux43cux435ux440-5.-ux43dux435ux43fux440ux435ux440ux44bux432ux43dux430ux44f-ux444ux443ux43dux43aux446ux438ux44f-ux43eux442-ux441ux43bux443ux447ux430ux439ux43dux43eux439-ux432ux435ux43bux438ux447ux438ux43dux44b}

Если \(X\) - случайная величина, то

\[
e^X,\qquad X^2,\qquad \sin X,\qquad |X|
\]

тоже случайные величины, потому что соответствующие функции
\(\mathbb R\to\mathbb R\) непрерывны и, следовательно, борелевски
измеримы.

\subsubsection{Пример 6. Не всякая функция является случайной
величиной}\label{ux43fux440ux438ux43cux435ux440-6.-ux43dux435-ux432ux441ux44fux43aux430ux44f-ux444ux443ux43dux43aux446ux438ux44f-ux44fux432ux43bux44fux435ux442ux441ux44f-ux441ux43bux443ux447ux430ux439ux43dux43eux439-ux432ux435ux43bux438ux447ux438ux43dux43eux439}

Пусть \(\Omega=\mathbb R\) и \(\mathcal F=\{\varnothing,\mathbb R\}\) -
тривиальная \(\sigma\)-алгебра. Функция

\[
X(\omega)=\omega
\]

не является случайной величиной относительно этой \(\mathcal F\), потому
что

\[
\{X\le0\}=(-\infty,0]\notin\mathcal F.
\]

Проблема не в самой формуле \(X(\omega)=\omega\), а в слишком бедной
\(\sigma\)-алгебре событий.

\subsection{6. Доказательства и
обоснования}\label{ux434ux43eux43aux430ux437ux430ux442ux435ux43bux44cux441ux442ux432ux430-ux438-ux43eux431ux43eux441ux43dux43eux432ux430ux43dux438ux44f}

\textbf{Почему в определении стоят прообразы, а не образы.}

Вероятность задана на подмножествах \(\Omega\), то есть на событиях из
\(\mathcal F\). Если в пространстве значений есть множество
\(B\subset E\), то событие ``значение \(X\) лежит в \(B\)'' находится в
исходном пространстве:

\[
\{\omega:X(\omega)\in B\}=X^{-1}(B).
\]

Именно его вероятность можно пытаться вычислить. Образ множества
\(X(A)\) лежит в \(E\) и сам по себе не является событием в \(\Omega\).

\textbf{Почему нельзя проверять только одно множество.}

Если проверить только \(\{X\le0\}\), это ничего не говорит о событиях
\(\{X\le x\}\) при других \(x\). Для распределения и интеграла нужно,
чтобы были измеримы все события вида \(\{X\in B\}\), где \(B\)
борелевское. Поэтому проверяют либо все борелевские множества, либо
порождающий класс, например все лучи \((-\infty,x]\).

\textbf{Почему координатный критерий работает только для правильной
\(\sigma\)-алгебры.}

Для \(\mathbb R^d\) с борелевской \(\sigma\)-алгеброй прямоугольники
порождают все борелевские множества, поэтому измеримость координат
достаточна. Если на пространстве значений взять произвольную более
богатую \(\sigma\)-алгебру, то координат может быть недостаточно: нужно
проверять именно \(\sigma\)-алгебру, выбранную в пространстве значений.

\textbf{Почему в метрическом пространстве часто используют счетную
базу.}

Если метрическое пространство сепарабельно, то у него есть счетная база
открытых шаров. Тогда для проверки борелевской измеримости достаточно
проверять прообразы шаров из этой базы. Это удобно для пространств
последовательностей и функций. Без сепарабельности некоторые привычные
упрощения могут не работать.

\textbf{Предупреждение про неограниченные и расширенные функции.}

Обычная вещественная случайная величина принимает значения в
\(\mathbb R\). Иногда разрешают значения \(\pm\infty\) и говорят о
расширенной случайной величине со значениями в

\[
\overline{\mathbb R}=[-\infty,\infty].
\]

Тогда нужно явно понимать, какая \(\sigma\)-алгебра стоит на
\(\overline{\mathbb R}\), и отдельно следить за корректностью операций
вроде \(+\infty-\infty\).

\textbf{Предупреждение про процессы с несчетным временем.}

Для счетного набора времен измеримость всех координат хорошо согласуется
с произведенной \(\sigma\)-алгеброй. Для процессов с несчетным
множеством времен, например \(t\in[0,1]\), утверждения о траекториях как
элементах пространства функций требуют дополнительных условий и выбора
конкретной \(\sigma\)-алгебры или топологии на пространстве траекторий.

\subsection{7. Что написать на
доске}\label{ux447ux442ux43e-ux43dux430ux43fux438ux441ux430ux442ux44c-ux43dux430-ux434ux43eux441ux43aux435}

\begin{enumerate}
\def\labelenumi{\arabic{enumi}.}
\tightlist
\item
  Измеримое отображение:
\end{enumerate}

\[
X:(\Omega,\mathcal F)\to(E,\mathcal E),
\qquad
X^{-1}(B)\in\mathcal F\quad \forall B\in\mathcal E.
\]

\begin{enumerate}
\def\labelenumi{\arabic{enumi}.}
\setcounter{enumi}{1}
\tightlist
\item
  Случайная величина:
\end{enumerate}

\[
X:(\Omega,\mathcal F,P)\to(\mathbb R,\mathcal B(\mathbb R)).
\]

\begin{enumerate}
\def\labelenumi{\arabic{enumi}.}
\setcounter{enumi}{2}
\tightlist
\item
  Критерий для \(\mathbb R\):
\end{enumerate}

\[
X\text{ измерима}
\quad\Longleftrightarrow\quad
\{X\le x\}\in\mathcal F\quad \forall x\in\mathbb R.
\]

\begin{enumerate}
\def\labelenumi{\arabic{enumi}.}
\setcounter{enumi}{3}
\tightlist
\item
  Проверка по порождающему классу:
\end{enumerate}

\[
\mathcal E=\sigma(\mathcal C),\quad
X^{-1}(C)\in\mathcal F\ \forall C\in\mathcal C
\Rightarrow
X\text{ измеримо}.
\]

\begin{enumerate}
\def\labelenumi{\arabic{enumi}.}
\setcounter{enumi}{4}
\tightlist
\item
  Векторный критерий:
\end{enumerate}

\[
X=(X_1,\ldots,X_d)\text{ измерим}
\quad\Longleftrightarrow\quad
X_1,\ldots,X_d\text{ измеримы}.
\]

\begin{enumerate}
\def\labelenumi{\arabic{enumi}.}
\setcounter{enumi}{5}
\tightlist
\item
  Композиция:
\end{enumerate}

\[
X\text{ измеримо},\quad g\text{ борелевски измерима}
\Rightarrow
g(X)\text{ измеримо}.
\]

\begin{enumerate}
\def\labelenumi{\arabic{enumi}.}
\setcounter{enumi}{6}
\tightlist
\item
  Случайная последовательность:
\end{enumerate}

\[
X:\Omega\to\mathbb R^{\mathbb N},
\qquad
X(\omega)=(X_1(\omega),X_2(\omega),\ldots).
\]

\subsection{8. Типичные дополнительные
вопросы}\label{ux442ux438ux43fux438ux447ux43dux44bux435-ux434ux43eux43fux43eux43bux43dux438ux442ux435ux43bux44cux43dux44bux435-ux432ux43eux43fux440ux43eux441ux44b}

\textbf{Вопрос. Зачем нужна измеримость случайной величины?}

Чтобы события вида \(\{X\in B\}\) принадлежали \(\mathcal F\) и имели
вероятность. Без измеримости нельзя корректно определить распределение
\(P_X(B)=P\{X\in B\}\).

\textbf{Вопрос. Почему в определении измеримости используется прообраз?}

Потому что вероятность задана на событиях в \(\Omega\). Множество \(B\)
лежит в пространстве значений, а событие в исходном пространстве равно
\(X^{-1}(B)\).

\textbf{Вопрос. Нужно ли проверять все борелевские множества на прямой?}

По определению да, но практически нет. Достаточно проверить порождающий
класс, например все лучи \((-\infty,x]\):

\[
\{X\le x\}\in\mathcal F\quad \forall x.
\]

\textbf{Вопрос. Почему непрерывная функция от случайной величины
измерима?}

Непрерывная функция борелевски измерима, а композиция измеримых
отображений измерима.

\textbf{Вопрос. Когда вектор \((X_1,\ldots,X_d)\) является случайным?}

Тогда и только тогда, когда все координаты \(X_1,\ldots,X_d\) являются
случайными величинами.

\textbf{Вопрос. Что такое случайный элемент метрического пространства?}

Это измеримое отображение из вероятностного пространства в метрическое
пространство с его борелевской \(\sigma\)-алгеброй:

\[
X:(\Omega,\mathcal F,P)\to(E,\mathcal B(E)).
\]

\textbf{Вопрос. Чем случайный элемент отличается от случайной величины?}

Случайная величина принимает значения в \(\mathbb R\). Случайный элемент
может принимать значения в более общем пространстве: \(\mathbb R^d\),
\(\mathbb R^{\mathbb N}\), пространстве функций, множестве состояний
цепи.

\textbf{Вопрос. Что такое траектория случайного процесса?}

При фиксированном \(\omega\) это функция времени:

\[
t\mapsto X_t(\omega).
\]

А при фиксированном \(t\) функция \(\omega\mapsto X_t(\omega)\)
называется сечением процесса.

\textbf{Вопрос. Может ли функция быть неслучайной из-за выбора
\(\sigma\)-алгебры?}

Да. Если \(\mathcal F\) слишком бедная, прообразы борелевских множеств
могут не лежать в \(\mathcal F\). Например, тождественная функция на
\(\mathbb R\) не измерима относительно тривиальной \(\sigma\)-алгебры
\(\{\varnothing,\mathbb R\}\).

\subsection{9. Типичные
ошибки}\label{ux442ux438ux43fux438ux447ux43dux44bux435-ux43eux448ux438ux431ux43aux438}

\begin{enumerate}
\def\labelenumi{\arabic{enumi}.}
\item
  Говорить ``случайная величина - любая функция на \(\Omega\)''.
  Правильно: случайная величина - измеримая функция.
\item
  Путать \(X(\omega)\) и распределение \(P_X\). Первое - значение
  функции при исходе, второе - мера на пространстве значений.
\item
  Проверять образ \(X(A)\) вместо прообраза \(X^{-1}(B)\). В определении
  измеримости всегда стоят прообразы.
\item
  Проверять одно событие, например \(\{X>0\}\), и считать, что
  измеримость доказана. Нужно проверить порождающий класс или всю
  \(\sigma\)-алгебру.
\item
  Забывать, какая \(\sigma\)-алгебра стоит в пространстве значений. Для
  вещественных величин обычно это \(\mathcal B(\mathbb R)\).
\item
  Считать, что координатная измеримость автоматически решает любую
  задачу со случайными функциями. Для пространств траекторий нужно
  учитывать выбранную \(\sigma\)-алгебру и топологию.
\item
  Путать сечение и траекторию процесса: \(X_t(\omega)\) как функция
  \(\omega\) - сечение, как функция \(t\) при фиксированном \(\omega\) -
  траектория.
\end{enumerate}

\subsection{10. Связи с другими
билетами}\label{ux441ux432ux44fux437ux438-ux441-ux434ux440ux443ux433ux438ux43cux438-ux431ux438ux43bux435ux442ux430ux43cux438}

\textbf{Билет 01.} Измеримость использует \(\sigma\)-алгебры как классы
допустимых событий.

\textbf{Билет 02.} Вероятностное пространство \((\Omega,\mathcal F,P)\)
- исходная область, на которой определяются случайные величины.

\textbf{Билет 04.} Для вещественных и векторных случайных величин
пространство значений снабжается борелевской \(\sigma\)-алгеброй.

\textbf{Билет 06.} Распределение случайного элемента определяется через
измеримость:

\[
P_X(B)=P(X^{-1}(B)).
\]

\textbf{Билет 07.} Интеграл Лебега строится сначала для измеримых
функций.

\textbf{Билет 08.} Характеристическая функция \(\mathbb E e^{itX}\)
имеет смысл, потому что \(e^{itX}\) измерима.

\textbf{Билет 10.} Условное математическое ожидание является измеримой
функцией относительно заданной \(\sigma\)-алгебры.

\textbf{Билет 15.} Теорема Колмогорова строит случайный процесс как
случайный элемент пространства последовательностей по согласованным
конечномерным распределениям.

\textbf{Билет 16.} Дискретный случайный процесс можно понимать как
последовательность случайных величин и как один случайный элемент
пространства последовательностей.

\subsection{11. Минимум на
удовлетворительно}\label{ux43cux438ux43dux438ux43cux443ux43c-ux43dux430-ux443ux434ux43eux432ux43bux435ux442ux432ux43eux440ux438ux442ux435ux43bux44cux43dux43e}

Нужно уметь сказать:

\begin{enumerate}
\def\labelenumi{\arabic{enumi}.}
\tightlist
\item
  Измеримое отображение:
\end{enumerate}

\[
X^{-1}(B)\in\mathcal F\quad \forall B\in\mathcal E.
\]

\begin{enumerate}
\def\labelenumi{\arabic{enumi}.}
\setcounter{enumi}{1}
\tightlist
\item
  Случайная величина - измеримое отображение
\end{enumerate}

\[
X:(\Omega,\mathcal F,P)\to(\mathbb R,\mathcal B(\mathbb R)).
\]

\begin{enumerate}
\def\labelenumi{\arabic{enumi}.}
\setcounter{enumi}{2}
\tightlist
\item
  Для вещественной случайной величины достаточно проверять
\end{enumerate}

\[
\{X\le x\}\in\mathcal F\quad \forall x.
\]

\begin{enumerate}
\def\labelenumi{\arabic{enumi}.}
\setcounter{enumi}{3}
\item
  Случайный вектор измерим тогда и только тогда, когда измеримы его
  координаты.
\item
  Один пример: \(\mathbf 1_A\) является случайной величиной, если
  \(A\in\mathcal F\).
\end{enumerate}

\subsection{12. Минимум на
хорошо}\label{ux43cux438ux43dux438ux43cux443ux43c-ux43dux430-ux445ux43eux440ux43eux448ux43e}

Кроме минимума на удовлетворительно, нужно:

\begin{enumerate}
\def\labelenumi{\arabic{enumi}.}
\tightlist
\item
  Объяснить, почему достаточно проверять порождающий класс.
\item
  Доказать или аккуратно проговорить критерий для \(\{X\le x\}\).
\item
  Сформулировать измеримость композиции.
\item
  Объяснить, почему непрерывная функция от случайной величины измерима.
\item
  Разобрать случайный вектор и случайную последовательность.
\end{enumerate}

\subsection{13. Что добавить для
отлично}\label{ux447ux442ux43e-ux434ux43eux431ux430ux432ux438ux442ux44c-ux434ux43bux44f-ux43eux442ux43bux438ux447ux43dux43e}

Для сильного ответа стоит добавить:

\begin{enumerate}
\def\labelenumi{\arabic{enumi}.}
\tightlist
\item
  Доказательство через класс
\end{enumerate}

\[
\mathcal D=\{B:X^{-1}(B)\in\mathcal F\}.
\]

\begin{enumerate}
\def\labelenumi{\arabic{enumi}.}
\setcounter{enumi}{1}
\tightlist
\item
  Координатный критерий для \(\mathbb R^d\) с доказательством через
  прямоугольники.
\item
  Связь случайной последовательности с произведенной \(\sigma\)-алгеброй
  и цилиндрами.
\item
  Различие между сечением и траекторией случайного процесса.
\item
  Предупреждение: для несчетного времени и пространств траекторий нужны
  дополнительные условия.
\end{enumerate}

\subsection{14. Устная версия ответа на 5
минут}\label{ux443ux441ux442ux43dux430ux44f-ux432ux435ux440ux441ux438ux44f-ux43eux442ux432ux435ux442ux430-ux43dux430-5-ux43cux438ux43dux443ux442}

Измеримые функции нужны для того, чтобы функции от исхода были
корректными вероятностными объектами. Пусть есть измеримые пространства
\((\Omega,\mathcal F)\) и \((E,\mathcal E)\). Отображение

\[
X:\Omega\to E
\]

называется измеримым, если

\[
X^{-1}(B)\in\mathcal F
\qquad \forall B\in\mathcal E.
\]

Это означает, что для любого наблюдаемого множества значений \(B\)
событие \(\{X\in B\}\) действительно является событием, которому можно
приписать вероятность.

Вещественная случайная величина - это измеримое отображение

\[
X:(\Omega,\mathcal F,P)\to(\mathbb R,\mathcal B(\mathbb R)).
\]

По определению нужно проверять все борелевские множества, но практически
достаточно проверять лучи:

\[
\{X\le x\}\in\mathcal F
\qquad \forall x\in\mathbb R.
\]

Это верно потому, что лучи \((-\infty,x]\) порождают
\(\mathcal B(\mathbb R)\), а класс множеств \(B\), для которых
\(X^{-1}(B)\) измерим, является \(\sigma\)-алгеброй.

Случайный вектор

\[
X=(X_1,\ldots,X_d)
\]

измерим тогда и только тогда, когда измеримы все координаты.
Доказательство идет через прямоугольники:

\[
\{X_1\le x_1,\ldots,X_d\le x_d\}
=
\bigcap_{k=1}^d\{X_k\le x_k\}.
\]

Если \(g\) - борелевски измеримая функция, а \(X\) - случайный элемент,
то \(g(X)\) снова измерим. В частности, непрерывные функции от случайных
величин, суммы, произведения, максимумы и минимумы случайных величин
снова являются случайными величинами.

Случайный элемент метрического пространства - это измеримое отображение
в \((E,\mathcal B(E))\). Например, случайная последовательность может
рассматриваться как одно отображение

\[
X:\Omega\to\mathbb R^{\mathbb N},
\qquad
X(\omega)=(X_1(\omega),X_2(\omega),\ldots),
\]

если на пространстве последовательностей взять произведенную
\(\sigma\)-алгебру. Это связывает билет с конечномерными
распределениями, теоремой Колмогорова и дискретными случайными
процессами.

\subsection{15. Если осталось 2 минуты перед
экзаменом}\label{ux435ux441ux43bux438-ux43eux441ux442ux430ux43bux43eux441ux44c-2-ux43cux438ux43dux443ux442ux44b-ux43fux435ux440ux435ux434-ux44dux43aux437ux430ux43cux435ux43dux43eux43c}

Запомнить:

\[
X:(\Omega,\mathcal F)\to(E,\mathcal E)
\text{ измеримо}
\quad\Longleftrightarrow\quad
X^{-1}(B)\in\mathcal F\ \forall B\in\mathcal E.
\]

Случайная величина:

\[
X:(\Omega,\mathcal F,P)\to(\mathbb R,\mathcal B(\mathbb R)).
\]

Для \(\mathbb R\):

\[
X\text{ измерима}
\quad\Longleftrightarrow\quad
\{X\le x\}\in\mathcal F\ \forall x.
\]

Для вектора:

\[
(X_1,\ldots,X_d)\text{ измерим}
\quad\Longleftrightarrow\quad
X_1,\ldots,X_d\text{ измеримы}.
\]

Главная мысль: измеримость нужна, чтобы \(\{X\in B\}\) было событием и
имело вероятность. Главная ловушка: случайная величина - не любая
функция, а измеримая функция.

\newpage

\end{document}
